\subsection{紧区间上函数列极限的积分}

把 II, p.~52 的定理 1 应用于紧区间上规则函数这一特殊情形，用积分的适当记号翻译出来就是：

命题 1. 设 A 是被一个滤子 \(\mathfrak{F}\) 所滤化的集合，\((\mathbf{f}_{\alpha})_{\alpha \in \mathrm{A}}\) 是紧区间 \(\mathrm{I} = [a, b]\) 上的一个规则函数族；若函数 \(\mathbf{f}_{\alpha}\) 在 I 上一致收敛于一个（规则的）函数 \(\mathbf{f}\)（关于滤子 \(\mathfrak{F}\)），则

\[
\lim _ {\mathfrak {F}} \int_ {a} ^ {b} \mathbf {f} _ {\alpha} (t) d t = \int_ {a} ^ {b} \mathbf {f} (t) d t.\tag{1}
\]

这个命题的两个推论在应用中很重要：

推论 1. 设 \((\mathbf{f}_n)\) 是紧区间 \(\mathrm{I} = [a, b]\) 上的一个规则函数序列。若序列 \((\mathbf{f}_n)\) 在 \(\mathrm{I}\) 上一致收敛于一个（规则的）函数 \(\mathbf{f}\)，则有

\[
\lim _ {n \rightarrow \infty} \int_ {a} ^ {b} \mathbf {f} _ {n} (t) d t = \int_ {a} ^ {b} \mathbf {f} (t) d t.\tag{2}
\]

特别地，若一个级数，其通项 \(u_{n}\) 是 I 上的规则函数，在 I 上一致收敛于 f，则以 \(\int_{a}^{b}\mathbf{u}_{n}(t)dt\) 为通项的级数收敛且其和为 \(\int_{a}^{b}\mathbf{f}(t)dt\)（“一致收敛级数的逐项积分”）。

推论 2. 设 A 是拓扑空间 F 的一个子集，f 是 I × A 到 R 上的完备赋范空间 E 中的一个映射，使得对每个 \(\alpha \in A\)，函数 \(x \mapsto \mathbf{f}(x, \alpha)\) 在 I 上是规则的。若函数 \(x \mapsto \mathbf{f}(x, \alpha)\) 在 I 上一致收敛于一个（规则的）函数 \(x \mapsto \mathbf{g}(x)\)（当 \(\alpha\) 趋于一点 \(\alpha_{0} \in \overline{A}\) 而保持属于 A 时），则有

\[
\lim _ {\alpha \rightarrow \alpha_ {0}, \alpha \in \mathrm{A}} \int_ {a} ^ {b} \mathbf {f} (x, \alpha) d x = \int_ {a} ^ {b} \mathbf {g} (x) d x.\tag{3}
\]

特别地：

命题 2（“积分对参数的连续性”）。设 F 是紧空间，I = {[}a, b{]} 是 R 中的紧区间，f 是 I × F 到 R 上的完备赋范空间 E 中的连续映射；则函数 \(\mathbf{h}(\alpha) = \int_{a}^{b} \mathbf{f}(x, \alpha) dx\) 在 F 上连续。

事实上，由于 \(\mathbf{f}\) 在紧空间 \(\mathrm{I} \times \mathrm{F}\) 上一致连续，故函数 \(\mathbf{f}(x, \alpha)\) 一致收敛于 \(\mathbf{f}(x, \alpha_0)\)（在 \(\mathrm{I}\) 上，当 \(\alpha\) 趋于任意一点 \(\alpha_0 \in \mathrm{F}\) 时）。

本命题的一个应用如下：函数 \((x, \alpha) \mapsto x^{\alpha}\) 在乘积 I × J 上连续，其中 I = {[}a, b{]} 是满足 0 \textless{} a \textless{} b 的紧区间，而 J 是 R 中的任意紧区间；由此断定 \(\int_{a}^{b} x^{\alpha} dx\) 是 \(\alpha\) 在 R 上的连续函数；而当 \(\alpha\) 为有理数且 \(\neq -1\) 时，这个函数等于 \(\frac{b^{\alpha+1} - a^{\alpha+1}}{\alpha + 1}\)，且函数 \(\alpha \mapsto \frac{b^{\alpha+1} - a^{\alpha+1}}{\alpha + 1}\) 在 R 的每个不含 -1 的区间上连续；于是（恒等式的延拓）有 \(\int_{a}^{b} x^{\alpha} dx = \frac{b^{\alpha+1} - a^{\alpha+1}}{\alpha + 1}\)（对一切实数 \(\alpha \neq -1\)）；这又意味着，对一切实数 \(\alpha\)，\(x^{\alpha}\) 的导数是 \(\alpha x^{\alpha-1}\)（参见 III, p.~94）。

